% `tex.web` §977 and §675: the box a `\vsplit` hands back is packed, and a
% packed box is reported on.
%
%     q:=vpackage(q,h,exactly,split_max_depth);
%
% `vpackage` is `vpack` with a depth limit, and §675 complains about the result
% exactly as it does for a `\vbox`: too high is `Overfull \vbox (Npt too high)`
% and too short is `Underfull \vbox (badness B)`, each `detected at line N` —
% or `has occurred while \output is active`, which is where trip.tex's is.
%
% We packed the split-off part in silence. It is the one packing path in this
% engine that said nothing, and the only reason it went unnoticed is that a
% `\vsplit` is usually of a box built to fit.
\catcode`\{=1 \catcode`\}=2
\showboxbreadth=2 \showboxdepth=1
\vbadness=0 \vfuzz=0pt \splitmaxdepth=2pt
\font\rip=cmr10 \rip

% Too short for what it was split to: underfull.
\setbox1\vbox{\hbox{}\vskip 100pt\hbox{}}
\setbox2\vsplit1 to 5pt
\message{PROBE:one=split}

% Too tall: overfull. The split takes everything down to the break, so a
% single tall box has to come whole.
\setbox3\vbox{\vbox to 40pt{}\penalty-10000 \hbox{}}
\setbox4\vsplit3 to 5pt
\message{PROBE:two=split}

% ...and one that fits says nothing.
\setbox5\vbox{\hbox to 0pt{}\penalty-10000 \hbox{}}
\setbox6\vsplit5 to 0pt
\message{PROBE:three=done}
\end
